%%%%%%%%%%%%%%%%%%%%%%%%%%%%%%%%%%%%%%%%%%%%%%%%%%%%%%%%%%%%%%%%
\appendix
\section{附录}
%%%%%%%%%%%%%%%%%%%%%%%%%%%%%%%%%%%%%%%%%%%%%%%%%%%%%%%%%%%%%%%%

\noindent\textbf{附录A：实验使用的Holt-Laury量表。}
收益域和损失域的10道配对彩票题目详见随报告提交的“附录/附录A\_实验材料/实验材料与问题手册.md”。\par
\noindent\textbf{附录B：实验条件的完整提示词。}
六组条件提示词分别为gain\_self、gain\_agent\_invisible、gain\_agent\_visible、loss\_self、loss\_agent\_invisible和loss\_agent\_visible，系统提示词另存于system文件，完整文本见随报告提交的“附录/附录B\_完整提示词”文件夹。\par
\noindent\textbf{附录C：分析代码与数据文件。}
所有统计分析的Python脚本、图表和中间数据文件均已分类保存在随报告提交的“附录/附录C\_数据代码与结果”目录下，包括scripts、data和figures三个子文件夹。
其中，修订后新增的收益域H1置换检验与H2差中差置换检验由“附录/附录C\_数据代码与结果/scripts/revision\_stats.py”生成，对应结果文件为“附录/附录C\_数据代码与结果/data/revision\_statistical\_tests.md”。
早期探索性统计摘要仅用于分析过程追踪，最终报告以正文表格和revision\_statistical\_tests.md为准。
